\newpage
\section{Related Work}\label{relatedwork}
The foundational work on post-quantum WireGuard was presented by Hülsing, Ning, Schwabe, Weber, and Zimmermann at IEEE S\&P 2021 \cite{pqwg}. Unlike most earlier work on post-quantum security for real-world protocols, their variant does not only address post-quantum confidentiality, but also targets full post-quantum authentication. To achieve this, they replace the Diffie-Hellman-based handshake with a more generic construction using only KEMs, resulting in a protocol that provides post-quantum security for both key exchange and identity authentication. Security is established in both the symbolic model, covering identity hiding and denial-of-service (DoS) mitigation among other properties, and the computational model of Dowling and Paterson \cite{dowling2018wireguard}, in which they prove key indistinguishability.

The core design challenge they address is the MTU constraint described in Section \ref{wg-PQC}: because WireGuard operates over UDP and assumes that handshake messages fit within a single unfragmented IPv6 packet, the key material contributed by PQC algorithms must fit within 1232 bytes. To satisfy this constraint, they instantiate their construction using Classic McEliece as the CCA-secure KEM and a passively secure variant of Saber,\footnote{CCA security guarantees indistinguishability against an adversary with access to a decryption oracle, while passive security (IND-CPA) only guarantees indistinguishability against an eavesdropper.} carefully selected so that the combined key material fits within the available space. A significant practical limitation of this choice is the large public key size of Classic McEliece, which substantially increases the server's memory footprint and complicates deployment in resource-constrained environments such as the Linux kernel.

PQ-WireGuard serves as an important performance reference point as the algorithms it uses, Classic McEliece and Saber, were the best available candidates at the time, but have since been superseded by the NIST-standardized ML-KEM and the newly selected HQC. This thesis revisits the performance question using these standardized algorithms, allowing for a more relevant comparison against the current state of post-quantum standardization.

Rosenpass \cite{rosenpass}, developed by Varner, Lipp, Schmidt, Dua, and Zaeske, takes a different approach to the integration problem. Rather than modifying the WireGuard handshake directly, Rosenpass runs as a separate process alongside WireGuard and performs a post-quantum key exchange independently. The resulting shared secret is then supplied to WireGuard as its pre-shared symmetric key (PSK), producing a hybrid connection where security holds as long as either the classical WireGuard handshake or the Rosenpass key exchange remains secure.

Such a sidecar architecture avoids any modification to WireGuard's codebase, making it practical to deploy today without kernel changes. The protocol builds on the PQ-WireGuard design but improves it by adding resistance against state disruption attacks, where an adversary replays the first handshake message to disrupt the responder's state. Rosenpass addresses this by returning an encrypted cookie to the initiator, so that the responder does not commit state until the initiator proves reachability. The implementation is written in Rust and uses Classic McEliece for long-term keys and CRYSTALS-Kyber-512 (Round 3 NIST) for ephemeral keys.

One important limitation of the PSK-injection strategy used by Rosenpass is that it only achieves post-quantum confidentiality as mentioned in Section \ref{sec:psk}. While Rosenpass itself provides post-quantum authentication via Classic McEliece, WireGuard's own static authentication keys remain Curve25519-based. Thus, a quantum adversary running Shor's algorithm could still break authentication independently of Rosenpass. Achieving full post-quantum authentication would therefore require replacing the static key mechanism entirely, which the PSK slot is not designed to do.

For this thesis, Rosenpass is relevant as a deployed reference implementation that uses CRYSTALS-Kyber in a production context. Its PSK-based integration strategy also represents a natural performance baseline. Since it adds a separate handshake on top of WireGuard, the combined latency is higher than a fully integrated solution. The performance results in this thesis can therefore be interpreted in the context of whether a deeper integration actually yields a meaningful latency improvement over the Rosenpass approach.

The most recent related work is by Hashimoto, Katsumata, Niot, and Wiggers, to appear at IEEE S\&P 2026 \cite{rebar}, which adresses the original PQ-WireGuard design by Hülsing et al. \cite{pqwg}, and improves it across three dimensions: protocol design, computational security, and efficiency.

The main design contribution of the paper is the introduction of \textit{reinforced KEMs} (RKEMs), a principled way to replace DH operations within the Noise framework with KEM operations. As KEMs produce two distinct values, a public key and a ciphertext, compared to a DH exchange that only produces a single shared value, the question of which values to mix into the key derivation chain, and in what order, arises.
 
The original PQ-WireGuard handled this substitution in an ad-hoc manner which leads to two consequences. First, it leaves the protocol vulnerable to unknown key share (UKS) attacks, where an adversary can exploit the absence of proper key binding to cause two parties to derive the same session key while disagreeing on each other's identity. This undermines the authentication guarantees of the protocol, as the two parties no longer agree on each other's identity despite holding the same session key.
Second, it causes improper PSK handling: since KEMs lack client-side authentication in the first round, Hülsing et al.\ were forced to absorb the PSK into the key derivation chain twice, a workaround with no clean cryptographic justification. The RKEM abstraction formalizes this by providing a precise interface that governs the way KEM outputs are absorbed into the handshake state.

The paper extends the computational security analysis beyond what was covered in the original work by Hülsing et al., which only proved key indistinguishability in the computational model and relied on the symbolic model for the remaining properties. The new work formalizes these additional security notions such as entity authentication, denial-of-service security, and identity hiding, directly in the computational model. It also fixes a subtle issue in the prior computational model of both Hülsing et al. and Dowling and Paterson \cite{dowling2018wireguard}, arising from the mid-protocol handling of pre-shared keys in the post-specified peer model. They do this by explicitly splitting the responder algorithm into an identification phase and a key derivation phase, so that the pre-shared key is only assigned once the initiator's identity has been established.

The most practically significant contribution of the RKEM-based redesign is that it eliminates the need for Classic McEliece as the client's long-term static key. By replacing this with Rebar, a novel ML-KEM-inspired RKEM construction, the server's per-client public key size is reduced by approximately 190 to 390$\times$, while still achieving full post-quantum authentication for both parties. The design also exploits WireGuard's out-of-band static key provisioning, described in Section \ref{sec:wg-keys}, allowing the handshake to remain within the UDP packet size limit.

The paper represents the current state of the art for formally analyzed post-quantum WireGuard designs, and it uses ML-KEM as its primary building block, the same algorithm that is evaluated in this thesis. While this thesis does not implement the RKEM construction, the performance characteristics of ML-KEM measured here are directly relevant to understanding the practical cost of the approach proposed by Hashimoto et al.

\section{State of the art}
The NIST Post-Quantum Cryptography Standardization Process \cite{nistpqc} continues in parallel with this thesis. ML-KEM was finalized as FIPS 203 in August 2024 and is the algorithm evaluated in this work. HQC was selected as a backup standard in March 2025 and is currently undergoing draft standardization, with a final standard expected in 2027 \cite{Moody2025NISTIR8545}. The ongoing standardization of HQC is relevant as the performance results for HQC reported in this thesis reflect the current reference implementation, not a final and optimized standard.

Post-quantum key exchange is actively being deployed in the Transport Layer Security (TLS) protocol, which represents the most widely monitored early indicator of PQC adoption at scale. As part of this effort, the IETF has published a draft standard for hybrid key exchange in TLS 1.3 using \texttt{X25519MLKEM768}, which combines X25519 and ML-KEM-768 into a single key share \cite{ietf-tls-ecdhe-mlkem-04}. 

Cloudflare began deploying \texttt{X25519MLKEM768} in TLS connections in early 2024 and published an analysis of its performance and adoption trajectory \cite{Westerbaan2025PQ}. Their findings indicate that hybrid ML-KEM adds minimal latency overhead for typical HTTPS connections, largely because TLS key material is transmitted in-band and the computational cost of ML-KEM is low. This supports the premise of this thesis that ML-KEM is viable in latency-sensitive protocols, though WireGuard imposes stricter constraints than TLS due to its MTU-bound message model. The TLS deployment data also provides an external reference against which the handshake latency results of this thesis can be contextualized.